\section{Experiments}
\label{experiments}

\subsection{Research Questions}
\label{exp:rqs}

The evaluation is organized around the following research questions:
\begin{itemize}[leftmargin=*]
  \item \textbf{RQ1:} How accurately does the workflow detect infeasible or conditionally feasible requirements?
  \item \textbf{RQ2:} How accurately does it identify the underlying constraint category?
  \item \textbf{RQ3:} How much do knowledge guidance and multi-agent role separation contribute to the results?
  \item \textbf{RQ4:} How much human verification effort can be reduced without removing human decision authority?
\end{itemize}

\subsection{Experimental Setup}
\label{exp:setup}

The experimental setup will define the input requirement format, available project knowledge, agent prompts, constraint rule sources, and human annotation protocol.
All implementation details should be reported sufficiently for reproduction.

\subsection{Datasets}
\label{exp:datasets}

The dataset should contain natural language requirements annotated with feasibility labels and constraint categories.
Candidate sources include course project requirements, public requirements engineering artifacts, and manually constructed infeasible variants.
The final paper should report dataset size, domains, annotation guidelines, and inter-annotator agreement.

\subsection{Baselines}
\label{exp:baselines}

Baselines should compare the proposed workflow against simpler feasibility review strategies:
\begin{itemize}[leftmargin=*]
  \item Single-agent requirement feasibility review.
  \item Rule-only constraint checker.
  \item Multi-agent review without knowledge guidance.
  \item Human-only review used as the reference process.
\end{itemize}

\subsection{Evaluation Metrics}
\label{exp:metrics}

\subsubsection{Feasibility Detection Accuracy}
\label{exp:metric_feasibility}

Measures whether the workflow correctly classifies requirements as feasible, conditionally feasible, infeasible, or insufficient information.

\subsubsection{Constraint Verification Accuracy}
\label{exp:metric_constraint}

Measures whether the identified constraint category matches the human-annotated reason for feasibility risk.

\subsubsection{Unrealistic Requirement Detection}
\label{exp:metric_unrealistic}

Measures precision, recall, and F1 for requirements judged unrealistic under the available project context.

\subsubsection{Human Verification Reduction}
\label{exp:metric_human}

Measures the reduction in manual review effort, such as fewer requirements requiring full expert inspection or fewer clarification rounds.

\subsection{Main Results}
\label{exp:main_results}

This subsection will report the main quantitative comparison after the dataset and baselines are finalized.
No experimental results are claimed at the structure-design stage.

\subsection{Ablation Study}
\label{exp:ablation}

The ablation study should isolate the effect of role separation, knowledge guidance, constraint rules, and human feedback.
The expected output is an evidence-based explanation of which components improve feasibility verification.

\subsection{Case Study}
\label{exp:case_study}

The case study should present a small number of representative requirements, agent findings, human judgments, and final refinements.
It should illustrate the feasibility analysis workflow without overgeneralizing from anecdotal examples.
